% Generated draft: requires mathematical and source review.
\documentclass[11pt]{article}
\usepackage[margin=25mm]{geometry}
\usepackage{fontspec}
\setmainfont{Latin Modern Roman}
\usepackage{amsmath,amssymb,mathtools,unicode-math}
\setmathfont{Latin Modern Math}
\usepackage{enumitem,needspace,xurl,hyperref}
\setlength{\parindent}{0pt}
\setlength{\parskip}{5pt}
\setlength{\emergencystretch}{3em}
\newcommand{\N}{\mathbb{N}}
\newcommand{\Z}{\mathbb{Z}}
\newcommand{\Q}{\mathbb{Q}}
\newcommand{\R}{\mathbb{R}}
\newcommand{\C}{\mathbb{C}}
\newcommand{\F}{\mathbb{F}}
\newcommand{\A}{\mathbb{A}}
\newcommand{\PP}{\mathbb{P}}
\newcommand{\E}{\mathbb{E}}
\newcommand{\eps}{\varepsilon}
\newcommand{\dd}{\,\mathrm{d}}
\begin{document}
\section*{UnsolvedMath ID: 10400093}
\textbf{Source problem code:} catalogue entry
\subsection{Definitions and mathematical statement}
Let \(T_n\) denote the trivial \(n\)-component link. For an unoriented skein quadruple \((L_0,L_1,L_2,L_3)\), \(L_{i+1}\) is obtained from \(L_i\) by a right-handed half twist on the same two arcs, as specified in the absent Figure 15. The conjecture, attributed to J. Przytycki, has four parts:
\begin{enumerate}
\item There exists an invariant of unoriented links \(P_1(L)\in\Z[x,t]\) with \(P_1(T_n)=t^n\) and
\[P_1(L_0)+xP_1(L_1)-xP_1(L_2)-P_1(L_3)=0.\]
\item There exists an invariant of unoriented framed links \(P_2(L)\in\Z[A,A^{-1},t]\) with \(P_2(T_n)=t^n\), \(P_2(L^{(1)})=-A^3P_2(L)\), and
\[P_2(L_0)+A(A^2+A^{-2})P_2(L_1)+(A^2+A^{-2})P_2(L_2)+AP_2(L_3)=0.\]
Here \(L^{(1)}\) is the source's positive half twist of the framing; this terminology is retained rather than replaced by a full-twist convention.
\item Put \(D=x+y+xy+y^2\). There exists an invariant of unoriented framed links
\[P_3(L)\in\Z[a,a^{-1},x,y,D^{-1}]\]
with \(P_3(L^{(1)})=aP_3(L)\),
\[P_3(L_0)+axP_3(L_1)+a^2yP_3(L_2)-a^3(x+y+1)P_3(L_3)=0,\]
and \(P_3(T_n)=c^{\,n-1}\). The flattened displayed initial condition appears to intend
\[c=\frac{-a^3(x+y+xy+x^2)+a^7(x+y+1)^2-a^{-1}}{x+y+xy+y^2},\]
but the fraction structure is lost in extraction; this reading requires source confirmation, with the numerator's \(x^2\) and denominator's \(y^2\) retained exactly as extracted.
\item For each of parts (1), (2), (3), the invariant, if it exists, is not unique: there are at least two distinct invariant functions satisfying that part's conditions. The source notes that all such functions agree on trivial links, hence on their span in the corresponding cubic skein module.
\end{enumerate}
The source states that (3) specializes to (1) at \(a=1,y=-x\), and to (2) at \(a=-A^3,x=-1-A^{-4},y=A^{-4}+A^{-8}\). These are recorded as proposed specialization claims. In the first, \(D\) becomes zero, so substitution is not a homomorphism from the displayed localization; cancellation or a specified extension is required. The parameter \(t\) in the first two targets is independent in those formulations, and the precise compatibility of these target rings under specialization is not supplied.

The preceding conjectures assume \(b_\infty=0\). The source then considers \(b_\infty\) invertible in a coefficient ring \(R\), with \(a\) invertible and skein parameters \(b_0,b_1,b_2,b_3,b_\infty\). Write \(t^n\) for \(T_n\), and \(H\) for the unoriented Hopf link. It reports the denominator relation
\[\bigl(a^{-3}b_3+a^{-2}b_2+a^{-1}b_1+b_0+b_\infty t\bigr)L=0\tag{33}\]
and the numerator and mirror relations
\begin{align*}
b_3(L\mathbin{\#}H)+(ab_2+b_1t+a^{-1}b_0+ab_\infty)L&=0,\\
b_0(L\mathbin{\#}H)+(a^{-1}b_1+b_2t+ab_3+a^2b_\infty)L&=0.
\end{align*}
Using the stated amphicheirality of the unoriented Hopf link, it obtains
\[\begin{split}
\bigl((b_0b_1-b_2b_3)t+(a^{-1}b_0^2-ab_3^2)
 +(ab_0b_2-a^{-1}b_1b_3)\\
 {}+b_\infty(ab_0-a^2b_3)\bigr)L=0.\tag{34}
\end{split}\]
The further question is whether the natural trivial-link map identifies
\[\mathcal S_{4,\infty}(S^3;R)\cong R[t]/I,\]
where \(I\) is generated by the left sides of (33) and (34) with \(L=t\), so that \(t^i\) represents the trivial \(i\)-component link. Its precise module and multiplication conventions depend on the missing surrounding definitions and Figure 16. Finally the source proposes \(b_0=b_3=a=1\), \(b_1=b_2=x\), \(b_\infty=y\), possibly leading to an invariant in \(\Z[x,y]\) satisfying
\[L_3+xL_2+xL_1+L_0+yL_\infty=0.\]
This is a proposed further construction; the local \(L_\infty\) diagram, normalization and relation to the preceding invertibility assumption are not supplied.
\subsection{Short English statement}
The conjecture asks for three cubic-skein link invariants with prescribed values and framing laws, and asserts their nonuniqueness. It also asks whether two displayed relations completely describe a related skein module of \(S^3\), followed by another proposed polynomial specialization.
\subsection{Sources}
\begin{itemize}
\item Local UnsolvedMath registry, permanent ID \(10400093\), exact source code \texttt{catalogue entry}.
\item J. Przytycki, source attribution, ``see [286]''; the full bibliographic entry for [286] is not supplied.
\item Ohtsuki, \emph{Problems on invariants of knots and 3-manifolds} (2002), Conjecture 4.15, PDF page 79; Figures 15--16 and equations (33)--(34): \url{https://msp.org/gtm/2002/04/gtm-2002-04-024s.pdf}. Registry access date 2026-07-29, pdf-font-structure extraction.
\end{itemize}
\end{document}
